\section{Experimental Evaluation Plan}\label{sec:evaluation}
All experiments in this section are prospective. No measurement table, usability improvement, or agent success rate has been obtained for \sys{}. Table~\ref{tab:eval} separates the intended evidence from the research questions. Protocol decisions and actual calibration parameters must be recorded before collecting test data.

\begin{table}[t]
\caption{Prospective evaluation matrix; every result is pending.}\label{tab:eval}
\centering\footnotesize\setlength{\tabcolsep}{4pt}
\begin{tabular}{@{}P{.08\columnwidth}P{.39\columnwidth}P{.44\columnwidth}@{}}\toprule
ID & Reference / comparison & Main measures\\\midrule
E1 & Known masses; independent reference balance & MAE, RMSE, maximum error, drift\\
E2 & Labeled physical actions; delta vs FSM vs confirmation & Precision, recall, F1, prompts\\
E3 & Scripted three-user, two-scale switches & Wrong attribution, success, latency\\
E4 & Matched tasks: manual vs semi-automatic & Time, actions, omissions, completion\\
E5 & Frozen tool fixtures and expected answers & Tool choice, values, sources, access safety\\\bottomrule
\end{tabular}
\end{table}

\subsection{E1: Weight Accuracy and Stability}
For each scale, reference loads spanning its documented operating range are placed repeatedly in randomized order. Record cold-start behavior, zero drift, settling time, and repeatability. Calibration and evaluation placements are kept separate. With errors $e_i=\hat m_i-m_i^*$ against an independent reference mass,
\begin{equation}\operatorname{MAE}=\frac{1}{N}\sum_{i=1}^N|e_i|,\quad
\operatorname{RMSE}=\sqrt{\frac{1}{N}\sum_{i=1}^Ne_i^2}.\label{eq:errors}\end{equation}
Propose five load levels $\{0,.2,.4,.6,.8\}M$ within documented capacity, five randomized placements per level in each of two sessions per scale: 50 test placements per scale. Calibration uses three separate placements at each of three reference levels (zero, $.4M$, $.8M$) and is fixed before that session's test block. Matching load values are allowed, but test placements are never used to fit coefficients. Record the exact physical masses, reference uncertainty, mounting and session conditions before collection. These are proposed counts, not completed trials. Report session/scale summaries and maximum absolute error as well. Reference-balance resolution and calibration uncertainty bound what can be inferred about the sensor. Test both assemblies rather than pooling them into an apparent increase in sample size.

\subsection{E2: Event Detection and Confirmation}
A labeled protocol includes drinking, lift-and-return without intake, refill, discard, spill, vessel change, and accidental contact. Ground truth combines a scripted action label, observer timing, and independently weighed transferred liquid. Ground-truth drink labels must not be produced by the same confirmation interface being evaluated. Compare (A) simple stable weight difference, (B) temporal FSM, and (C) FSM with human confirmation on the same observation traces. Retain all exclusions and timeouts.

Propose ten trials for each of the seven action classes, spread over two acquisition blocks, giving 70 action opportunities. Pilot/calibration traces are excluded. Freeze matching tolerance $\delta_t$ before test scoring using observer timing resolution and a justified timing allowance. Match predicted events to ground-truth ingestion events within $\delta_t$ of the labeled return using maximum-cardinality one-to-one matching; minimize total absolute time difference to break ties, then use event IDs. A matched event is TP, every unmatched recorded-ingestion prediction is FP (including duplicate detections), and every unmatched ground-truth drink is FN. Refills and spills are negative action opportunities; their false drink records count as FP. Unresolved/time-out drinks remain FN for final-record scoring, and unresolved coverage is reported separately. Also report specificity over labeled negative opportunities.

Evaluate raw candidates and final records separately; ordinary inferred records count as final predictions unless rejected. Confirmation uses actual logged user responses, not oracle labels injected into the algorithm. A replay without such responses evaluates only the sensor stages. All methods see identical traces and fixed scoring rules. For matched events, report
\begin{equation}P=\frac{TP}{TP+FP},\quad R=\frac{TP}{TP+FN},\quad
F_1=\frac{2PR}{P+R}.\label{eq:detection}\end{equation}
Zero-denominator quantities are reported as undefined, not silently perfect. Report raw candidate performance separately from final recorded-intake performance. Count prompts, user corrections, and unresolved events. A high precision obtained by asking about every event must remain visible as interaction burden.

\subsection{E3: Shared-Device Attribution}
Three test users alternate control of both scales according to a scripted schedule. Include switches during an absent-cup interval, food weighing, and queued offline events. An external schedule establishes the intended event owner. Report wrong-owner records divided by attributable events, switch acknowledgement success, elapsed switch time, and unresolved events. Propose six scenarios (ordinary switch, mid-lift switch, duplicate replay, lost acknowledgement, disconnection and restart), each for three intended owners and both scales: 36 scripted attempts. The expected outcomes follow the switch barrier in Section~\ref{sec:method}, including holds rather than guessed owners. Alongside wrong-owner/resolved-record rate, report resolved coverage over all eligible events and wrong-owner/all-eligible rate. Report activated switches over all requested switches; failed attempts and missing acknowledgements remain in the denominator. Trace event IDs and session versions across retries. This is a consistency evaluation, not a measurement of identity recognition.

\subsection{E4: Recording Effort}
Use a within-person crossover pilot with the three intended users performing matched hydration and food tasks in manual and semi-automatic conditions. Plan two sessions per participant, each with four matched tasks per condition (two hydration, two food). Use AB in one session and BA in the other, assigning the initial order before testing and rotating task variants: 16 planned tasks per person. Log actual order and deviations; training tasks are separate. Provide equivalent training and task quantities. Measure task time, explicit actions, omissions relative to the observer log, and completed tasks. Include confirmation and correction time, rather than counting only automatic records. Set a proposed 10-minute task deadline before collection. Retain timed-out tasks as incomplete with time reported as censored at the deadline; do not turn those into fast successful trials. Report each person's completion/omission proportions and median paired time differences for completed matched tasks, alongside counts of excluded incomplete pairs. Do not impute completion times or pool repeated tasks as independent people. Report paired participant-level summaries; this small convenience sample supports feasibility observations, not population-level effectiveness claims.

\subsection{E5: Agent Reliability and Authorization}
Propose 30 frozen cases: five each for daily hydration, multi-day hydration, nutrition, compound queries, missing/conflicting data, and authorization/adversarial requests. Run each case three times from a reset fixture and conversation, for 90 planned runs. Freeze model/provider identifier, sampling settings, prompt, tool schema, workflow version and fixture hash. Include day-boundary queries, a corrected source revision, malicious source text, forged identities, unavailable tools and upstream failures. These are test stimuli, not observations of a running agent.

A deterministic oracle computes quantities from the fixture's authorized current revisions and accepted unit conversions. Freeze acceptable tool-call sequences (including equivalent permitted alternatives), rounding rules and expected statuses before runs. Score tool correctness $T_{ik}$, numeric/status correctness $N_{ik}$, source/version correctness $P_{ik}$ and authorization/safety $A_{ik}$ as binary components for case $i$, repeat $k$. Numeric tolerance is half of the declared last displayed unit after canonical conversion; no-data cases require the expected nonnumeric status, not a fabricated zero. An unauthorized call attempt fails $A_{ik}$ even if the backend denies it; backend non-disclosure is reported as a separate defense check. Strict run success is
\begin{equation}S_{ik}=T_{ik}N_{ik}P_{ik}A_{ik},\quad
\bar S=\frac{1}{90}\sum_{i=1}^{30}\sum_{k=1}^{3}S_{ik}.\label{eq:agent}\end{equation}
Report component rates and family-level counts as well as empirical all-three-repeat case reliability,
\begin{equation}C_3=\frac{1}{30}\sum_{i=1}^{30}\prod_{k=1}^{3}S_{ik}.\label{eq:repeat}\end{equation}
This read-only metric is inspired by repeated evaluation~\cite{yao2024tau}, not a reproduction of its state-changing benchmark. Grade final answers and returned sources against the oracle, not by another unconstrained LLM. Have a second evaluator check semantic/adversarial judgments and record disagreements. Refusals and missing-data statements must match the case's expected policy. A correct number cannot compensate for wrong provenance or a privacy failure. No clinical safety guarantee follows from this small scripted suite.

\subsection{Reproducibility and Data Handling}
Before execution, freeze sensor parameters, software versions, stimulus schedules, case definitions and analysis rules. Keep acquisition-block and case identifiers for clustered descriptive summaries; do not treat temporally adjacent sensor windows or repeated agent runs as independent participants. Planned counts are feasibility choices, not a power calculation. Save de-identified traces with record and correction relationships. Participant consent and applicable institutional review must precede real personal-data collection. Raw health records and identifiable screenshots do not belong in the public release. Disclose missing trials, device failures, exclusions and changes to the protocol.
